% SPDX-License-Identifier: GPL-3.0-or-later
% Copyright (C) 2026 Alessandro Greco
\chapter{Истоки: mix-сети, ремейлеры, Crowds и Freedom}\label{cap:02-radici}
\aggiornatoal{23 сентября 2026 г.}{исторические источники, цитируемые в этой главе}

\begin{obiettivi}
К концу этой главы вы сможете:
\begin{itemize}[nosep]
  \item объяснять идею \emph{mix}, предложенную Дэвидом Чаумом, и её значение для систем анонимности;
  \item различать высоколатентные системы для электронной почты и низколатентные системы для веб-сёрфинга;
  \item проследить развитие ремейлеров от anon.penet.fi до Mixminion;
  \item определить место Crowds, PipeNet и Freedom в истории, приведшей к Tor.
\end{itemize}
\end{obiettivi}

Tor возник не на пустом месте. На сегодняшний день это очередной этап более чем сорокалетней истории исследований и реально развёрнутых систем. Некоторые из них добились успеха, многие оказались неудачными, но каждая чему-то нас научила. В этой главе прослеживаются основные этапы развития вплоть до появления onion routing.

% ---------------------------------------------------------------------------
\section{Mix Чаума (1981)}\label{sec:radici-chaum}

В 1981 году Дэвид Чаум опубликовал пятистраничную статью, положившую начало целому направлению исследований: \titolo{Untraceable Electronic Mail, Return Addresses, and Digital Pseudonyms}~\cite{chaum1981untraceable}. Задача касалась электронной почты: как отправить сообщение так, чтобы наблюдатель не мог связать отправителя с получателем, даже если вся сеть находится под наблюдением?

Решением стал \textbf{mix}: сервер, который принимает сообщения, зашифрованные его открытым ключом, расшифровывает их, собирает в группу (\emph{batch}) и пересылает вместе в порядке, отличном от порядка поступления. Связать входящее сообщение с исходящим мешают три приёма:
\begin{itemize}
  \item исходящее сообщение выглядит совершенно иначе, чем соответствующее входящее, поскольку mix снял один слой шифрования;
  \item все сообщения имеют одинаковую длину;
  \item сообщения перемешиваются и выпускаются вместе, поэтому даже время передачи не связывает их друг с другом.
\end{itemize}
\begin{figure}[htbp]
\centering
% SPDX-License-Identifier: GPL-3.0-or-later
% Copyright (C) 2026 Alessandro Greco
%% Русское альтернативное описание EPUB.
\altfig{Микс-сервер Чаума принимает слева четыре зашифрованных сообщения одинаковой длины с метками A, B, C и D. Собрав группу, сняв слой шифрования и перемешав сообщения, он выпускает справа четыре одинаково оформленных сообщения с независимыми метками 1–4. Схема не раскрывает соответствие между входами и выходами.}
% Микс Чаума: сбор группы, снятие одного слоя и перемешивание.
\begin{tikzpicture}[font=\sffamily\scriptsize,
  msg/.style={draw,minimum width=1.5cm,minimum height=0.5cm,rounded corners=1pt}]
  \node[draw=torviola,fill=torviola!8,rounded corners=4pt,minimum width=2.4cm,minimum height=3.4cm] (mix) at (5.2,0) {};
  \node[font=\sffamily\small,text=torviola] at (5.2,1.35) {\textbf{mix}};
  \node[align=center,text=torviola] at (5.2,0.4) {собрать группу\\снять слой\\перемешать};
  \foreach \i/\y/\c in {A/1.2/torverde!30,B/0.4/torviola!25,C/-0.4/avviso!25,D/-1.2/torgrigio!25} {
    \node[msg,fill=\c] (in\i) at (1.6,\y) {\i};
    \node[font=\sffamily\tiny,left=1pt of in\i] {отпр.\ \i};
  }
  \foreach \i/\y in {1/1.2,2/0.4,3/-0.4,4/-1.2} {
    \node[msg,fill=torgrigio!20] (out\i) at (8.8,\y) {\i};
  }
  \foreach \i in {A,B,C,D} \draw[->,torgrigio] (in\i) -- (mix.west |- in\i);
  \foreach \i in {1,2,3,4} \draw[->,torgrigio] (mix.east |- out\i) -- (out\i);
  \node[align=center,font=\sffamily\tiny] at (1.6,-1.9) {одинаковая длина,\\зашифровано для mix};
  \node[align=center,font=\sffamily\tiny] at (8.8,-1.9) {другой вид,\\перемешанный порядок};
  \draw[decorate,decoration={brace,amplitude=4pt}] (0.6,1.55) -- node[above=4pt] {входящая группа} (2.5,1.55);
  \draw[decorate,decoration={brace,amplitude=4pt}] (7.9,1.55) -- node[above=4pt] {выходящая группа} (9.7,1.55);
\end{tikzpicture}

\caption{Mix Чаума. Наблюдатель, видящий одновременно вход и выход, не может связать сообщения: все они имеют одинаковую длину, их внешний вид меняется после снятия mix-сервером слоя шифрования, а выходят они одновременно и в другом порядке.}
\label{fig:mix}
\end{figure}

Вместо доверия одному mix отправитель может пропустить сообщение через \textbf{цепочку} mix-серверов, послойно зашифровав его ключом каждого из них. Каждый mix снимает свой слой и узнаёт только следующий узел. Если хотя бы один mix в цепочке честен, связь между отправителем и получателем остаётся скрытой. Именно здесь находится зародыш идеи onion routing.

В той же статье Чаум предложил \emph{анонимные обратные адреса}, позволяющие получателям отвечать, не зная отправителя, а также цифровые псевдонимы. В 1988 году он добавил второй механизм --- \emph{DC-nets}~\cite{chaum1988dining}, обеспечивающий безусловную анонимность ценой очень высоких затрат и оказавший преимущественно теоретическое влияние.

% ---------------------------------------------------------------------------
\section{Высокая и низкая задержка}\label{sec:radici-latenza}

Mix Чаума обеспечивает сильную защиту, но работает \textbf{медленно}: сообщение должно ждать на каждом mix, пока сформируется пакет, и доставка может занять минуты или часы. Для электронной почты это приемлемо, но не для веб-сёрфинга, где страница должна загружаться за секунды.

Отсюда возникает фундаментальное разделение:
\begin{description}
  \item[Высоколатентные системы] (mix-сети, ремейлеры) задерживают и перемешивают сообщения и способны противостоять наблюдателю, видящему всю сеть, но подходят лишь для коммуникаций, допускающих задержки.
  \item[Низколатентные системы] (onion routing, Tor) пересылают данные сразу после поступления, позволяя просматривать веб-страницы, вести чаты и поддерживать интерактивные соединения; однако наблюдатель, видящий оба конца, может сопоставлять временные характеристики и объём трафика.
\end{description}
Tor сознательно выбирает второй путь: система, которой почти никто не пользуется из-за чрезмерной медлительности, никого не защищает (\cref{sec:anon-utenti}).

% ---------------------------------------------------------------------------
\section{Ремейлеры}\label{sec:radici-remailer}

В 1990-х идеи Чаума воплотились в реальных сервисах, называемых \textbf{ремейлерами}: серверах, которые получают электронное письмо и пересылают его после удаления сведений об отправителе.

\paragraph{anon.penet.fi (1993--1996).} Самый известный ремейлер работал в Финляндии под управлением Йохана \enquote{Julf} Хельсингиуса с 1993 по 1996 год. Это был не mix: сервис присваивал каждому пользователю \emph{псевдоним} и хранил таблицу соответствий между псевдонимами и реальными адресами, чтобы доставлять ответы. Со временем система обрабатывала тысячи сообщений в день. Но именно эта таблица в итоге привела к закрытию сервиса. В 1995 и 1996 годах Церковь саентологии при содействии финской полиции получила личности некоторых пользователей, распространявших документы церкви. В 1996 году Хельсингиус закрыл сервис, сославшись на правовую неопределённость вокруг Интернета в Финляндии. Урок, который впоследствии прямо учитывали разработчики onion routing~\cite{syverson2011peel}, заключается в том, что система, сохраняющая соответствие между пользователями и их личностями, создаёт единую точку отказа: оператора можно обязать раскрыть эти данные.

\paragraph{Cypherpunk, Mixmaster, Mixminion.} Тем временем сообщество \emph{cypherpunk} разрабатывало ремейлеры без таблиц соответствий: отправители послойно шифровали сообщения для цепочки серверов. Последующие поколения --- сначала Mixmaster, затем Mixminion~\cite{danezis2003mixminion} --- ввели сообщения фиксированной длины, настоящее пакетирование и защиту от повторной передачи сообщений. Роджер Динглдайн и Ник Мэтьюсон были среди авторов Mixminion, выпущенного в 2003 году; примерно в то же время они проектировали Tor.

% ---------------------------------------------------------------------------
\section{К низкой задержке}\label{sec:radici-bassa}

\paragraph{PipeNet (1995).} В феврале 1995 года Вэй Дай предложил в рассылке cypherpunk систему PipeNet --- низколатентную сеть, в которой все участники непрерывно передавали бы трафик с постоянной скоростью, не оставляя признаков для корреляции. Такая система сопротивлялась бы временному анализу, но неприемлемой ценой: если хотя бы один участник прекращал передачу, останавливаться должна была вся сеть~\cite{syverson2011peel}. PipeNet так и не был реализован, но остаётся ориентиром того, что потребовалось бы для противодействия корреляции трафика.

\paragraph{Onion routing (1996).} В том же году в Naval Research Laboratory формировалась концепция onion routing. Эту историю мы разберём в \cref{cap:03-nascita-tor}.

\paragraph{Crowds (1998).} Майкл Рейтер и Авиэль Рубин предложили Crowds~\cite{reiter1998crowds}. Пользователи образуют \enquote{толпу}, а каждый веб-запрос передаётся случайно выбранному участнику. С некоторой вероятностью тот отправляет запрос на веб-сайт, иначе передаёт следующему участнику. Веб-сайт видит запрос, пришедший от произвольного члена группы, а ни один участник не может определить, был ли человек, передавший ему запрос, его первоначальным источником. Crowds не использует послойное шифрование и защищает только от веб-сайта и других участников, но не от наблюдателя сети. Тем не менее система ввела строгий способ выражать анонимность через вероятность, а также анализ атак сговаривающихся участников, предвосхитивший исследования guard-реле.

\paragraph{Freedom (1999--2001).} Крупнейшей сетью анонимности до Tor была Freedom, которой управляла канадская компания Zero-Knowledge Systems. Её анонсировали в конце 1998 года, а работала она с конца 1999 до конца 2001 года~\cite{syverson2011peel}. Это был коммерческий сервис: пользователи платили подписку, а серверы обслуживали компания и её партнёры. Сервис закрылся, потому что доходов от подписок оказалось недостаточно для его поддержки. Его история показывает два урока, которым Tor на практике последовал, как мы увидим далее: сеть анонимности устойчивее, когда ею управляет множество независимых сторон, а не одна компания, и модель финансирования важна не меньше криптографии (\cref{sec:nascita-comunita}).

\begin{approfondimento}[Много «луковиц»]
В тот же период предлагалось множество других систем onion routing, включая полностью децентрализованные Tarzan и MorphMix, а также Cebolla, у которой Tor заимствовал идею строить цепь по одному переходу за раз, выполняя обмен Диффи--Хеллмана для каждого перехода~\cite{syverson2011peel}. Параллельно с Tor появились и другие проекты, пошедшие иными путями: I2P, ориентированный прежде всего на внутренние сервисы и действующий до сих пор, и AN.ON, созданный в Дрезденском университете вместе с клиентом JAP (и предлагавшийся также как коммерческий сервис под названием JonDonym), который использует фиксированные цепочки mix-серверов под управлением известных организаций. Каждый проект по-своему балансирует между децентрализацией, доверием и производительностью.
\end{approfondimento}

% ---------------------------------------------------------------------------
\section*{Итоги}
\begin{itemize}
  \item Mix Чаума (1981) скрывает связь между входящими и исходящими сообщениями, изменяя их внешний вид, выравнивая длину и перемешивая порядок. Цепочка mix-серверов с послойным шифрованием предвосхищает «луковицу».
  \item Высоколатентные системы противостоят сильным наблюдателям, но медленны; низколатентные пригодны для интерактивной работы, однако уязвимы к временной корреляции.
  \item anon.penet.fi продемонстрировал опасность хранения идентичностей пользователей; Mixmaster и Mixminion развили подход mix-сетей для электронной почты.
  \item PipeNet, Crowds и Freedom исследовали низкую задержку с разными компромиссами. Freedom показал важность сети, управляемой множеством сторон, и устойчивой модели финансирования.
\end{itemize}

% ---------------------------------------------------------------------------
\section*{Упражнения}

\begin{esercizio}[1]\label{es:02-mix}
Mix собирает пакеты по 100 сообщений. Наблюдатель видит, как все 100 сообщений входят и как все 100 выходят. Если наблюдателю известен отправитель одного входящего сообщения, какова вероятность правильно угадать соответствующее исходящее сообщение при случайном выборе? Что изменится, если сам наблюдатель отправил 99 из этих 100 сообщений?
\end{esercizio}

\begin{esercizio}[2]\label{es:02-penet}
Набросайте архитектуру ремейлера, позволяющего отвечать пользователям без хранения таблицы соответствий между псевдонимом и адресом. Какие идеи Чаума для этого понадобятся?
\end{esercizio}

\begin{esercizio}[2]\label{es:02-crowds}
В Crowds при вероятности пересылки $p_f = 3/4$ через сколько участников в среднем проходит запрос, прежде чем попасть на веб-сайт? Почему участник, получивший запрос, не может быть уверен, что передавший его участник и был первоначальным источником?
\end{esercizio}

% ---------------------------------------------------------------------------
\section*{Дополнительная литература}
\begin{itemize}
  \item D.~Chaum, \titolo{Untraceable Electronic Mail, Return Addresses, and Digital Pseudonyms}~\cite{chaum1981untraceable}.
  \item M.~Reiter, A.~Rubin, \titolo{Crowds: Anonymity for Web Transactions}~\cite{reiter1998crowds}.
  \item G.~Danezis, R.~Dingledine, N.~Mathewson, \titolo{Mixminion: Design of a Type III Anonymous Remailer Protocol}~\cite{danezis2003mixminion}.
  \item \emph{Anonymity Bibliography} проекта Free Haven, \url{https://www.freehaven.net/anonbib/}: справочная коллекция научных публикаций по этой теме.
\end{itemize}
